\subsection{FREE COMMUTATIVE GROUPS AND MONOIDS}

The set \(\mathbf{Z}^{\mathbf{x}}\) of all mappings of \(\mathbf{X}\) into \(\mathbf{Z}\) is a commutative group under the law defined by \((\alpha + \beta)(x) = \alpha(x) + \beta(x)(\alpha, \beta \in \mathbf{Z}^{\mathbf{x}}, x \in \mathbf{X})\) ; the elements of \(\mathbf{Z}^{\mathbf{x}}\) are sometimes called multiindices. The identity element, denoted by 0, is the constant mapping with value 0. For \(\alpha \in \mathbf{Z}^{\mathbf{x}}\) , the set \(S_{\alpha}\) of \(x \in \mathbf{X}\) such that \(\alpha(x) \neq 0\) is called the support of \(\alpha\) ; then \(S_0 = \varnothing\) and \(S_{\alpha + \beta} \subset S_{\alpha} \cup S_{\beta}\) for \(\alpha, \beta\) in \(\mathbf{Z}^{\mathbf{x}}\) . Therefore the set \(\mathbf{Z}^{(\mathbf{x})}\) of mappings \(\alpha: \mathbf{X} \to \mathbf{Z}\) of finite support is a subgroup of \(\mathbf{Z}^{\mathbf{x}}\) called the free commutative group constructed on \(\mathbf{X}\) .

For all \(x \in \mathbf{X}\) , let \(\delta_x\) denote the element of \(\mathbf{Z}^{(\mathbf{x})}\) defined by

\[
\delta_ {x} (y) = \left\{ \begin{array}{l l} 1 & \text { if } y = x \\ 0 & \text { if } y \neq x. \end{array} \right.\tag{10}
\]

Also, for \(\alpha \in \mathbf{Z}^{(\mathbf{x})}\) , the integer \(|\alpha|\) , the length of \(\alpha\) is defined by the formula

\[
| \alpha | = \sum_ {x \in X} \alpha (x).\tag{11}
\]

The following relations are immediately established:

\[
\alpha = \sum_ {x \in \mathbf {X}} \alpha (x) \cdot \delta_ {x}\tag{12}
\]

\[
\left| \delta_ {x} \right| = 1, \quad \left| 0 \right| = 0\tag{13}
\]

\[
| \alpha + \beta | = | \alpha | + | \beta |\tag{14}
\]

for \(\alpha, \beta\) in \(\mathbf{Z}^{(\mathbf{X})}\) and \(x\) in \(\mathbf{X}\) .

The order relation \(\alpha \leqslant \beta\) is defined in \(\mathbf{Z}^{(\mathbf{x})}\) by \(\alpha(x) \leqslant \beta(x)\) for all \(x \in \mathbf{X}\) . The relations \(\alpha \leqslant \beta\) and \(\alpha' \leqslant \beta'\) imply \(\alpha + \alpha' \leqslant \beta + \beta'\) , \(|\alpha| \leqslant |\beta|\) and \(-\alpha \geqslant -\beta\) ; further, the relation \(\alpha \leqslant \beta\) is equivalent to \(\beta - \alpha \geqslant 0\) . The set of elements \(\alpha \geqslant 0\) in \(\mathbf{Z}^{(\mathbf{x})}\) is denoted by \(\mathbf{N}^{(\mathbf{x})}\) ; it is the set of mappings of X into N of finite support and it is a submonoid of \(\mathbf{Z}^{(\mathbf{x})}\) called the free commutative monoid constructed on X. The elements of length 1 are the minimal elements in \(\mathbf{N}^{(\mathbf{x})} = \{0\}\) and constitute the set of \(\delta_x\) ( \(x \in \mathbf{X}\) ).

The monoid \(\mathbf{N}^{(\mathbf{x})}\) and the group \(\mathbf{Z}^{(\mathbf{x})}\) enjoy the following universal property.

PROPOSITION 10. Let M be a commutative monoid (resp. group) and f a mapping of X into M. There exists one and only one homomorphism of \(\mathbf{N}^{(\mathbf{x})}\) (resp. \(\mathbf{Z}^{(\mathbf{x})}\) ) into M such that \(g(\delta_x) = f(x)\) for all \(x \in \mathbf{X}\) . If \(\mathbf{M}\) is written additively, then \(g(\alpha) = \sum_{x \in \mathbf{X}} \alpha(x) \cdot f(x)\) for all \(\alpha\) in \(\mathbf{N}^{(\mathbf{x})}\) (resp. \(\mathbf{Z}^{(\mathbf{x})}\) ).

Let \(g\) be a homomorphism of \(\mathbf{N}^{(\mathbf{X})}\) (resp. \(\mathbf{Z}^{(\mathbf{X})}\) ) into \(\mathbf{M}\) such that \(g(\delta_x) = f(x)\) for all \(x \in \mathbf{X}\) . For all \(\alpha\) in \(\mathbf{N}^{(\mathbf{X})}\) (resp. \(\mathbf{Z}^{(\mathbf{X})}\) ), it follows from (12) that

\[
g (\alpha) = \sum_ {x \in \mathbf {X}} \alpha (x). g (\delta_ {x}) = \sum_ {x \in \mathbf {X}} \alpha (x). f (x),
\]

whence the uniqueness of g.

For all \(\alpha\) in \(\mathbf{N}^{(\mathbf{x})}\) (resp. \(\mathbf{Z}^{(\mathbf{x})}\) ) we write \(g(\alpha) = \sum_{x\in \mathbf{X}}\alpha (x).f(x)\) . Then obviously \(g(0) = 0\) ; for \(\alpha, \beta\) in \(\mathbf{N}^{(\mathbf{x})}\) (resp. \(\mathbf{Z}^{(\mathbf{x})}\) ),

\[
\begin{array}{r l} g (\alpha + \beta) & = \sum_ {x \in \mathbf {X}} (\alpha (x) + \beta (x)). f (x) \\ & = \sum_ {x \in \mathbf {X}} [ \alpha (x). f (x) + \beta (x). f (x) ] \\ & = \sum_ {x \in \mathbf {X}} \alpha (x). f (x) + \sum_ {x \in \mathbf {X}} \beta (x). f (x) \\ & = g (\alpha) + g (\beta) \end{array}
\]

and hence \(g\) is a homomorphism of \(\mathbf{N}^{(\mathbf{x})}\) (resp. \(\mathbf{Z}^{(\mathbf{x})}\) ) into \(\mathbf{M}\) . Also, for \(y\) in \(\mathbf{X}\) ,

\[
g (\delta_ {y}) = \sum_ {x \in \mathbf {X}} \delta_ {y} (x) \cdot f (x);
\]

now \(\delta_y(x).f(x) = 0\) for \(x\neq y\) and \(\delta_y(y).f(y) = f(y)\) , whence \(g(\delta_y) = f(y)\) .

Let \(u: X \to Y\) be a mapping. By Proposition 10 there exists one and only one homomorphism of \(\mathbf{Z}^{(X)}\) into \(\mathbf{Z}^{(Y)}\) which maps \(\delta_x\) to \(\delta_{u(x)}\) for all \(x \in X\) . It is denoted by \(\mathbf{Z}^{(w)}\) ; it is immediately seen that it maps \(\alpha \in \mathbf{Z}^{(X)}\) to the element \(\beta \in \mathbf{Z}^{(Y)}\) defined by

\[
\beta (y) = \sum_ {x \in u ^ {- 1} (y)} \alpha (x).\tag{15}
\]

As in the case of magmas (no. 1), the formula \(\mathbf{Z}^{(v\circ u)} = \mathbf{Z}^{(v)}\circ \mathbf{Z}^{(y)}\) is established for every mapping \(v\colon \mathbf{Y}\to \mathbf{Z}\) ; it is also shown that \(\mathbf{Z}^{(u)}\) is injective (resp. surjective, bijective) if \(u\) is.

Let \(S\) be a subset of \(X\) ; if \(i\) is the injection of \(S\) into \(X\) , the mapping \(f = \mathbf{Z}^{(i)}\) is an isomorphism of \(\mathbf{Z}^{(S)}\) onto the subgroup \(H\) of \(\mathbf{Z}^{(X)}\) generated by the elements \(\delta_s\) for \(s \in S\) . By (15),

\[
(f (\alpha)) (x) = \left\{ \begin{array}{l l} \alpha (x) & \text { if } x \in S \\ 0 & \text { if } x \in X - S \end{array} \right.
\]

and therefore H is the set of elements of \(\mathbf{Z}^{(\mathbf{x})}\) of support contained in S. Henceforth \(\mathbf{Z}^{(\mathbf{s})}\) will be identified with H by means of f.

Formula (15) shows that the restriction of \(\mathbf{Z}^{(u)}\) to \(\mathbf{N}^{(\mathbf{X})}\) induces a homomorphism \(\mathbf{N}^{(u)}\) of \(\mathbf{N}^{(\mathbf{X})}\) into \(\mathbf{N}^{(\mathbf{Y})}\) . Then \(\mathbf{N}^{(v\circ u)} = \mathbf{N}^{(v)}\circ \mathbf{N}^{(u)}\) for every mapping \(v\colon \mathbf{Y}\to \mathbf{Z}\) ; further, \(\mathbf{N}^{(u)}\) is injective (resp. surjective, bijective) if \(u\) is. If \(S\) is a subset of \(\mathbf{X}\) ,

\[
\mathbf {N} ^ {(\mathbf {S})} = \mathbf {Z} ^ {(\mathbf {S})} \cap \mathbf {N} ^ {(\mathbf {X})}.
\]

Remark. Let M be the multiplicative monoid of strictly positive integers and let P be the set of prime numbers (§ 4, no. 10, Definition 15). By Proposition 10 there exists a homomorphism u of \(\mathbf{N}^{(\mathfrak{P})}\) into M characterized by \(u(\delta_{p}) = p\) for every prime number p.~Then \(u(\alpha) = \prod_{p \in \mathfrak{P}} p^{\alpha(p)}\) for \(\alpha\) in \(\mathbf{N}^{(\mathfrak{P})}\) and Theorem 7 of § 4, no. 10 shows that u is an isomorphism of \(\mathbf{N}^{(\mathfrak{P})}\) onto H.

